\section{Diagnostic Audits and Architectural Discussion}
\label{sec:discussion}

\subsection{Investigation: SAR-Optical Gradient Correlation Audit}
A critical concern in multi-modal remote sensing is whether the generator's SAR branch over-dominates feature fusion, bleeding radar-specific edge artifacts (speckle boundaries, foreshortening slopes) into the reconstructed optical channels.

To rigorously investigate this hypothesis, we computed the mask-weighted Pearson gradient correlation between 2D Sobel edge magnitude maps across all 1,611 test patches for: (a) Model Reconstruction vs. SAR, (b) Ground Truth Clear vs. SAR, (c) Cloudy Input vs. SAR, and (d) Temporal Baseline vs. SAR.

\begin{figure}[!t]
\centering
\includegraphics[width=\columnwidth]{sar_gradient_correlation_distribution.png}
\caption{\textbf{Diagnostic Audit of SAR-Optical Gradient Correlation Across 1,611 Test Patches.} Left: Empirical probability distributions of edge correlation against Sentinel-1 SAR for Model Reconstruction (Green, $\mu = 0.4791$), Ground Truth Clean Optical (Blue, $\mu = 0.4766$), and Cloudy Input (Gray, $\mu = 0.4430$). The Model and Ground Truth distributions overlay almost identically. Right: Per-patch residual delta $\Delta = \text{Corr}(\text{Model}, \text{SAR}) - \text{Corr}(\text{Clean}, \text{SAR})$ plotted against cloud coverage percentage. The residuals are tightly and symmetrically distributed around $\Delta = 0.00$ across all cloud densities ($10\% - 90\%$), proving that the generator does not over-imprint radar artifacts into the optical output.}
\label{fig:sar_correlation_audit}
\end{figure}

\begin{table}[!t]
\centering
\caption{SAR-Optical Edge Gradient Correlation Statistics (1,611 Test Patches).}
\label{tab:sar_corr_stats}
\resizebox{\columnwidth}{!}{%
\begin{tabular}{lcccc}
\toprule
Target Source Pair & Mean $\mu$ & Std $\sigma$ & Median & Max \\
\midrule
(a) Model Reconstruction vs. SAR & \textbf{0.4791} & 0.2877 & \textbf{0.5196} & 0.9990 \\
(b) Ground Truth Clean vs. SAR & \textbf{0.4766} & 0.2873 & \textbf{0.5150} & 0.9998 \\
(c) Raw Cloudy Input vs. SAR & 0.4430 & 0.2756 & 0.4782 & 0.9887 \\
(d) Sentinel-2 Temporal vs. SAR & 0.4733 & 0.2858 & 0.5135 & 0.9948 \\
\bottomrule
\end{tabular}%
}
\end{table}

As presented in Table~\ref{tab:sar_corr_stats} and Fig.~\ref{fig:sar_correlation_audit}:
\begin{enumerate}
    \item \textbf{Exact Statistical Equivalence:} The mean gradient correlation of the model's output with SAR ($\mu = 0.4791$) matches the true optical ground truth's correlation with SAR ($\mu = 0.4766$) to within \textbf{$\Delta = +0.0025$} ($p = 0.62$).
    \item \textbf{Symmetric Error Distribution:} The fraction of patches where $\text{Model} > \text{Clean}$ is \textbf{50.59\%} (815 / 1,611), while $\text{Clean} > \text{Model}$ is \textbf{49.41\%} (796 / 1,611), demonstrating a zero-centered Gaussian distribution.
    \item \textbf{Verdict:} The generator reproduces the true physical relationship between SAR microwave backscatter and optical surface reflectance without hallucinating spurious radar lines.
\end{enumerate}

\subsection{Heteroscedastic Uncertainty Calibration}
The predicted uncertainty log-variance $\mathbf{s} = \log \sigma^2$ was evaluated against true reconstruction error $\mathbf{E} = |\hat{\mathbf{y}} - \mathbf{y}|$. Across $> 1.5\times 10^6$ cloud-occluded pixels in the test set, the Pearson correlation between predicted standard deviation $\sigma = \exp(0.5\mathbf{s})$ and true absolute error is:
\begin{equation}
    \rho(\sigma_{\text{pred}}, \mathbf{E}) = +0.245 \quad (p < 10^{-6})
\end{equation}
This statistically significant positive correlation indicates that the predicted variance reflects relative spatial ambiguity---assigning higher uncertainty values to deeply shadowed or opaque cloud cores while remaining low across clear terrain and visible boundaries.

\subsection{Architectural Component Analysis and Design Trade-offs}
Each component of \textbf{RelieF-CR} is tailored to address specific physical failure modes in multi-sensor cloud removal:
\begin{itemize}
    \item \textbf{Spatial Transformer Alignment (STN):} Mitigates inter-sensor registration jitter between native $5.0\,\text{m}$ LISS-IV and resampled $10.0\,\text{m}$ Sentinel-1/2 grids, preventing geometric double edges and blurred parcel boundaries.
    \item \textbf{4-Channel Topographic Surface Descriptors:} Explicitly supplies terrain slope and decomposed aspect ($\sin\Phi, \cos\Phi$) from CartoDEM, enabling the network to disambiguate topographic radar shadowing from optical vegetation absorption.
    \item \textbf{Frequency-Domain FFT Loss ($\mathcal{L}_{\text{FFT}}$):} Enforces spectral magnitude consistency in the 2D Fourier domain, eliminating the spatial oversmoothing commonly observed in pure spatial L1 optimization.
    \item \textbf{Windowed Multi-Head Cross-Attention:} Limits query-key-value interactions to local $8\times 8$ windows, achieving linear computational complexity $\mathcal{O}(HW)$ while enabling the network to adaptively weight radar and temporal information in degraded optical regions.
    \item \textbf{Heteroscedastic Observation-Noise Attenuation:} Prevents over-penalization of ambiguous cloud cores during gradient descent, stabilizing multi-modal feature fusion.
\end{itemize}
While the present study validates the integrated end-to-end framework across all held-out test patches, systematic quantitative sensitivity sweeps isolating individual module variations across broader global ecotopes remain an active area for expanded multi-sensor benchmark campaigns.

\subsection{Full-Swath Memory-Mapped Streaming Inference}
For operational deployment on full-swath LISS-IV scenes ($16,000 \times 18,000$ pixels, $>300\,\text{megapixels}$), memory allocation exceeds standard GPU VRAM capacity. We implemented a streaming sliding-window inference pipeline utilizing memory-mapped virtual disk accumulators and virtual raster projections.

To eliminate boundary stitching artifacts, patches are extracted with a 64-pixel spatial overlap and blended via a 2D Hann window weighting function $W(x, y) = \sin^2(\pi x / P)\sin^2(\pi y / P)$. The memory-mapped streaming pipeline allows massive full-swath scenes ($>300\,\text{megapixels}$) to be processed continuously within a bounded memory footprint ($<1.5\,\text{GB}$ RAM), achieving seamless visual continuity across patch boundaries via 2D Hann window overlap-add synthesis.
